% Options for packages loaded elsewhere
\PassOptionsToPackage{unicode}{hyperref}
\PassOptionsToPackage{hyphens}{url}
\documentclass[
]{ltjsarticle}
\usepackage{xcolor}
\usepackage{amsmath,amssymb}
\setcounter{secnumdepth}{-\maxdimen} % remove section numbering
\usepackage{iftex}
\ifPDFTeX
  \usepackage[T1]{fontenc}
  \usepackage[utf8]{inputenc}
  \usepackage{textcomp} % provide euro and other symbols
\else % if luatex or xetex
  \usepackage{unicode-math} % this also loads fontspec
  \defaultfontfeatures{Scale=MatchLowercase}
  \defaultfontfeatures[\rmfamily]{Ligatures=TeX,Scale=1}
\fi
\usepackage{lmodern}
\ifPDFTeX\else
  % xetex/luatex font selection
\fi
% Use upquote if available, for straight quotes in verbatim environments
\IfFileExists{upquote.sty}{\usepackage{upquote}}{}
\IfFileExists{microtype.sty}{% use microtype if available
  \usepackage[]{microtype}
  \UseMicrotypeSet[protrusion]{basicmath} % disable protrusion for tt fonts
}{}
\makeatletter
\@ifundefined{KOMAClassName}{% if non-KOMA class
  \IfFileExists{parskip.sty}{%
    \usepackage{parskip}
  }{% else
    \setlength{\parindent}{0pt}
    \setlength{\parskip}{6pt plus 2pt minus 1pt}}
}{% if KOMA class
  \KOMAoptions{parskip=half}}
\makeatother
\setlength{\emergencystretch}{3em} % prevent overfull lines
\providecommand{\tightlist}{%
  \setlength{\itemsep}{0pt}\setlength{\parskip}{0pt}}
\usepackage{newunicodechar}
\newunicodechar{→}{\ensuremath{\rightarrow}}
\newunicodechar{⟨}{\ensuremath{\langle}}
\newunicodechar{⟩}{\ensuremath{\rangle}}
\newunicodechar{⇒}{\ensuremath{\Rightarrow}}
\newunicodechar{⟹}{\ensuremath{\Longrightarrow}}
\newunicodechar{×}{\ensuremath{\times}}
\newunicodechar{≈}{\ensuremath{\approx}}
\newunicodechar{∼}{\ensuremath{\sim}}
\newunicodechar{≤}{\ensuremath{\le}}
\newunicodechar{≥}{\ensuremath{\ge}}
\newunicodechar{≠}{\ensuremath{\ne}}
\newunicodechar{∫}{\ensuremath{\int}}
\newunicodechar{−}{\ensuremath{-}}
\newunicodechar{✗}{\ensuremath{\times}}
\newunicodechar{✓}{\ensuremath{\checkmark}}
\usepackage{bookmark}
\IfFileExists{xurl.sty}{\usepackage{xurl}}{} % add URL line breaks if available
\urlstyle{same}
\hypersetup{
  hidelinks,
  pdfcreator={LaTeX via pandoc}}

\author{}
\date{}

\begin{document}

\section{A Survey of Related Prior Work for the Observer--System Beat
Conjecture --- Its Relation to Prior Programs and the Locus of the
Testable Regime
(Supplement)}\label{a-survey-of-related-prior-work-for-the-observersystem-beat-conjecture-its-relation-to-prior-programs-and-the-locus-of-the-testable-regime-supplement}

\textbf{Subtitle}: A record of a post-publication literature survey.
Related prior work that may be relevant is cited only where necessary;
the connection to the preceding paper is left as future work.

\textbf{Noriaki Kihara} (This is a \textbf{supplement} to the preceding
paper {[}1{]}. It is not a claim of established physics and contains no
new derivation or physical claim. After the publication of {[}1{]}, we
surveyed prior work that may be related to its mechanism, and this
records the result. Whether the cited works connect logically to the
conjecture of {[}1{]} is undetermined and is \textbf{future work}. This
supplement does not claim that they support the conjecture of {[}1{]};
it cites only the necessary parts within the range where relevance is
plausible, and organizes the possible relations.)

Version v0.1 (2026-06-28) DOI (Version): 10.5281/zenodo.20985750 (this
version, v0.1) DOI (Concept): 10.5281/zenodo.20985749 (for external
reference; auto-redirects to the latest version) Zenodo:
https://zenodo.org/records/20985750

\begin{center}\rule{0.5\linewidth}{0.5pt}\end{center}

\subsection{Abstract}\label{abstract}

The preceding paper {[}1{]} (hereafter ``the beat conjecture'')
proposed, as a candidate mechanism for the ``origin of randomness'' left
as a postulate by the localized-kernel paper {[}2{]}, that the
observer--system relative phase (beat/heterodyne) equidistributes (Weyl)
when the frequency ratio is irrational; combined with the
reproducing-kernel property, the shape of the position-resolved
intensity profile coincides with \(|\psi_{\rm base}|^2\) (the core),
distinguished from the transition from intensity to single clicks (the
bridge). In v0.4 the preceding paper does \textbf{not} claim the
finite-\(N\) \(O(1/N)\) correction as a measurable deviation; it
positions the beat mechanism as one facet of indeterminacy and adopts
the hypothesis that a more fundamental indeterminacy may prevent the
very realization of a fully coherent (strictly integer-ratio)
observation --- such a mechanism may exist (its existence is not
asserted) --- so that, if it does, the Born rule of standard physics
holds \textbf{more strongly} {[}1{]}.

This supplement records a literature survey carried out after the
publication of the beat conjecture. There are two conclusions. First,
\textbf{no prior work that proposes, tests, or observes the central
mechanism of the beat conjecture (the emergence of Born statistics from
the system--observer frequency beat plus irrational-ratio Weyl
equidistribution, and a Born-rule deviation in the coherent/commensurate
regime) in the same form was found within the scope surveyed}. Second,
there nonetheless exist several prior programs that may \textbf{share
the problem-consciousness} ``to derive the Born rule from a
deterministic substrate and threshold detection as an
emergent/approximate rule,'' and the beat conjecture may be related to
them. This supplement organizes these into (1) prior programs that may
share the problem-consciousness and (2) work of a different direction,
and additionally records that the precision-spectroscopy (atomic-clock)
domain that could test the beat conjecture has not, within the range
surveyed, performed the discriminating measurement (a measurement of
single-shot projection statistics as a function of the commensurability
of the relative phase). None of the citations claims identity of
mechanism.

\begin{center}\rule{0.5\linewidth}{0.5pt}\end{center}

\subsection{1. Introduction: motivation and place of this
supplement}\label{introduction-motivation-and-place-of-this-supplement}

The beat conjecture {[}1{]} is a conjecture (prediction) paper, not a
claim of established physics. After its publication, in order to locate
it accurately within the lineages of existing research, we surveyed
prior work that may be related.

This work has two purposes. First, to secure honesty: presenting a new
framework without grounding it in the existing literature would mislead
the scope of the claim. Second, to carve out the originality of the
conjecture (if any) accurately, as a difference from the known lineages.

This supplement merely records the result of the survey and does
\textbf{not} claim the following.

\begin{itemize}
\tightlist
\item
  It does not claim that the cited prior work supports the beat
  conjecture.
\item
  It does not claim that the cited prior work is the same mechanism as
  the beat conjecture.
\item
  It does not claim that any experiment has tested the beat conjecture.
\end{itemize}

Citations are limited to the necessary parts where relevance is
plausible. For works already cited by {[}1{]} (such as {[}6{]}{[}8{]}),
the content is re-checked in this supplement and the possible relation
to the conjecture is re-organized.

\begin{center}\rule{0.5\linewidth}{0.5pt}\end{center}

\subsection{2. Method and scope of the
survey}\label{method-and-scope-of-the-survey}

The survey centered on four areas: (i) foundational attempts to derive
the Born rule from a deterministic substrate, classical fields, and
threshold detection; (ii) tests of the Born rule using a coherent local
oscillator; (iii) the behavior of quantum projection noise (QPN) near
the lock point in atomic clocks and the presence/absence of
non-entanglement sub-QPN anomalies; and (iv) the application of the
method of arbitrary functions to the Born rule. Each candidate work was
checked against the primary source (published version or arXiv
original).

This survey is not exhaustive. A finite number of works were checked,
and the conclusion ``no such prior work exists'' is a statement about
the range surveyed, not a proof of absence across the whole literature.

\begin{center}\rule{0.5\linewidth}{0.5pt}\end{center}

\subsection{3. Prior programs that may be related (same
problem-consciousness)}\label{prior-programs-that-may-be-related-same-problem-consciousness}

There exist several programs that seek to derive the Born rule, not as a
basic postulate, but as an emergent/approximate rule from a deeper
deterministic layer and the detection process. These may \textbf{share
the problem-consciousness} with the beat conjecture. On the other hand,
whether the elements of the beat conjecture --- the beat,
irrational-ratio equidistribution, and a commensurability-dependent
deviation --- connect to these as a mechanism is, at present,
undetermined and is future work.

\subsubsection{3.1 Emergence of the Born rule from classical fields +
threshold detection (La
Cour--Williamson)}\label{emergence-of-the-born-rule-from-classical-fields-threshold-detection-la-courwilliamson}

La Cour and Williamson {[}3{]} (\emph{Quantum} 4, 350, 2020) showed, in
the context of quantum optics, that several phenomena thought to be
quantum can be reproduced using only classical fields, a deterministic
detection model of ``intensity threshold detectors,'' and a QED vacuum
assumed real (not virtual). The picture is that the quantum--classical
boundary is reproduced simply by applying standard conventions of data
analysis to discrete detection events.

The problem-consciousness (the emergence of the Born rule from a
deterministic substrate + threshold detection) may be close to the beat
conjecture. Whether the conjecture's logic of identifying the origin of
randomness with the observer--system frequency beat connects to this
work is not judged in this supplement and is left as future work.

\subsubsection{3.2 Prequantum classical statistical field theory (PCSFT,
Khrennikov)}\label{prequantum-classical-statistical-field-theory-pcsft-khrennikov}

Khrennikov's prequantum classical statistical field theory (PCSFT)
{[}4{]} (\emph{Phys. Lett. A} 372, 6588, 2008; related: {[}5{]}) regards
quantum particles as symbolic representations of ``prequantum random
fields,'' and claims that the probabilities arising when classical
random fields interact with threshold detectors agree with the Born rule
\textbf{only approximately}. The Born rule emerges from PCSFT as an
approximate theory, and as a consequence a \textbf{deviation from the
Born rule is predicted}.

PCSFT is a framework already cited by the beat conjecture {[}1{]} as
lineage A ({[}1{]} §6.1), and may be closest in spirit to this
conjecture in regarding ``the Born rule as an approximation to a
deterministic substrate, with a predicted deviation.'' Two points are
added. First, whether the deviation predicted by PCSFT connects to the
beat conjecture in a form dependent on the observer--system phase
coherence or commensurability is, at present, undetermined and is future
work. Second, since the early triple-slit experiments that Khrennikov
referenced in {[}4{]} as corroboration of the deviation, high-precision
Sorkin-type tests have detected no significant deviation from the Born
rule (see §4.1). That is, the central prediction of PCSFT is at present
not experimentally confirmed, and this supplement does not treat it as
supporting evidence.

\subsubsection{3.3 The method of arbitrary functions (Feintzeig et
al.)}\label{the-method-of-arbitrary-functions-feintzeig-et-al.}

Bonds, Burson, Feintzeig et al.~{[}6{]} (arXiv:2409.16457, 2024) applied
the ``method of arbitrary functions'' --- treating a dynamical parameter
as a random variable with an arbitrary initial distribution --- to
unitary Schrödinger dynamics, and showed that the Born-rule
probabilities appear as universal limiting behavior in the joint
long-time and small-\(\hbar\) limit. This is a framework cited by the
beat conjecture {[}1{]} as lineage B ({[}1{]} §4, §6.2).

The conjecture may be related to {[}6{]} in that it identifies the
``equidistributing dynamical parameter'' that {[}6{]} places abstractly
with the observer--system heterodyne relative phase. {[}6{]} itself,
within the range checked in this supplement, contains no experimental
proposal or prediction of a Born-rule deviation. Whether the two connect
as a mechanism is future work.

\subsubsection{3.4 Common points and differences
(summary)}\label{common-points-and-differences-summary}

The three programs of §3.1--3.3 may all share with the beat conjecture
the problem-consciousness ``to derive the Born rule from a
deterministic/classical substrate and the detection process as an
emergent/approximate rule.'' On the other hand, whether the elements of
the beat conjecture --- the observer--system frequency beat,
irrational-ratio Weyl equidistribution, and a qualitative break of the
Born rule in the coherent/commensurate regime --- connect to these as a
mechanism is not judged in this supplement and is left as future work.
The very fact that several independent routes may be heading toward the
same question, ``the Born rule may not be a basic postulate,'' is itself
worth recording.

\begin{center}\rule{0.5\linewidth}{0.5pt}\end{center}

\subsection{4. Possible relation to work of a different
direction}\label{possible-relation-to-work-of-a-different-direction}

\subsubsection{4.1 A Born-rule test using a coherent local oscillator
(Conlon et
al.)}\label{a-born-rule-test-using-a-coherent-local-oscillator-conlon-et-al.}

Conlon et al.~{[}7{]} (\emph{New J. Phys.} 26, 053003, 2024) tested the
Born rule (third-order / Sorkin interference) and the complexity of
Hilbert space using a three-arm interferometer with coherent states of
light and homodyne detection. The results --- Sorkin parameter
\(\kappa = 0.002 \pm 0.004\) and Peres parameter
\(F = 1.0000 \pm 0.0003\) --- are consistent with standard quantum
mechanics (no deviation from the Born rule was detected).

This is, within the range surveyed, one of the few Born-rule tests that
actually use a coherent local oscillator (homodyne), and may be related
to the beat conjecture. Its measured object is higher-order (Sorkin)
interference; whether it connects to the discriminant of the beat
conjecture {[}1{]} --- how the single-shot projection statistics at a
fixed observer--system relative phase depend on the commensurability
(rational/irrational) --- is not judged in this supplement and is future
work. We record that a Born-rule test using a coherent LO already
exists, and that, within the range surveyed, it does not include a
measurement in the discriminating regime of this conjecture.

\subsubsection{4.2 A deterministic substrate but a stance reluctant
about Born-rule deviation ('t
Hooft)}\label{a-deterministic-substrate-but-a-stance-reluctant-about-born-rule-deviation-t-hooft}

't Hooft's cellular-automaton interpretation {[}8{]} (Springer,
\emph{Fundamental Theories of Physics} 185, 2016) regards quantum
mechanics as a tool for treating a fundamental deterministic system.
This is an adjacent lineage already cited by the beat conjecture {[}1{]}
({[}1{]} §6.3).

't Hooft's framework, while placing a deterministic substrate, is said
to be reluctant about reproducible deviations from the Born rule (it
tends to keep the Born rule as a map of the relative probabilities of
ontic states). Given that the preceding paper v0.4 {[}1{]} adopts the
stance that ``a more fundamental indeterminacy may prevent a fully
coherent observation from being realized (the existence of such a
mechanism is a hypothesis, not asserted), and if so the Born rule of
standard physics holds more strongly,'' the beat conjecture (v0.4) and
't Hooft are --- for different reasons --- close in conclusion, in that
neither positively asserts an observable deviation from the Born rule.
PCSFT (§3.2), which predicts an observable deviation, differs in
direction from both on this point. We record that within the same
problem-consciousness of a ``deterministic substrate'' there are such
gradations.

\begin{center}\rule{0.5\linewidth}{0.5pt}\end{center}

\subsection{5. The locus of the testable regime: precision spectroscopy
has not performed the discriminating
measurement}\label{the-locus-of-the-testable-regime-precision-spectroscopy-has-not-performed-the-discriminating-measurement}

The beat conjecture {[}1{]} is operationally grounded in the Ramsey
metrology of atomic clocks (the relative phase
\(\varphi_0=(\nu_{\rm atom}-\nu_{\rm LO})\,T\) of the observer = local
oscillator and the atom, read by heterodyne). Hence atomic clocks are
the most natural testbed for this conjecture. This section records the
result of checking whether that literature performs the measurement
needed for the discrimination.

In the standard noise description of atomic clocks, quantum projection
noise (QPN) is treated as an \textbf{irreducible binomial noise floor}
arising from the probabilistic nature of measurement (Born-type
projection) {[}9{]} (Itano et al., \emph{Phys. Rev.~A} 47, 3554, 1993;
for a review, {[}10{]} Ludlow et al., \emph{Rev.~Mod. Phys.} 87, 637,
2015). In actual machines, no anomaly is reported in which sub-QPN /
non-binomial single-shot statistics appear near the lock point. Even in
configurations that partially raise the atom--LO coherence by long
interrogation or dead-time removal ({[}11{]} Schioppo et al.,
\emph{Nature Photonics} 11, 48, 2017; {[}12{]} Al-Masoudi et al.,
\emph{Phys. Rev.~A} 92, 063814, 2015), performance is governed by the
QPN floor and no anomaly is seen.

Demonstrations below the QPN / standard quantum limit are all
\textbf{deliberately created by spin squeezing (entanglement)} and
certified by entanglement witnesses such as the Wineland parameter
(e.g., {[}13{]} Malia et al., \emph{Phys. Rev.~Lett.} 125, 043202,
2020). This supplement treats these neither as corroboration nor as a
counter-example of the beat conjecture. The non-squeezing constraints of
atomic clocks (the Dick effect, laser phase noise, the coherence-time
limit) are also described as classical noise ({[}14{]} Schulte et al.,
\emph{Nat. Commun.} 11, 5955, 2020).

The point is this. Atomic clocks are routinely operated under the
conditions of on-resonance, phase-locked LO, and projective readout ---
just short of the discriminating regime of the beat conjecture --- but
the measurement needed for the discrimination (taking the statistics of
the single-shot excited fraction at a fixed atom--LO relative phase and
looking for non-binomial structure as a function of commensurability
(rational/irrational)) is, within the range of literature surveyed, not
performed in any of them. Normal operation is effectively in the
equidistributed regime (the irrational side) due to dither and noise,
where the predictions of the beat conjecture and of standard quantum
mechanics coincide. The fixed-coherence regime where the two diverge is,
being unnecessary for clock-making, not entered. This is not concealment
but a consequence of the measurement protocol structurally not entering
this regime.

The preceding paper v0.4 {[}1{]} does not claim this \(O(1/N)\)
correction as a measurable deviation; it adopts the hypothesis that a
more fundamental indeterminacy may prevent the very realization of a
fully coherent (strictly integer-ratio) observation --- such a mechanism
may exist (its existence is not asserted) --- so that, if it does, the
Born rule of standard physics holds more strongly. The fact this
supplement confirmed --- ``no anomaly is reported near the lock point /
the measurement needed for the discrimination is not performed'' ---
does not contradict this stance (the non-observation of the deviation is
consistent with the preceding paper's expectation). We record having
checked the presence/absence of the discriminating measurement.

\begin{center}\rule{0.5\linewidth}{0.5pt}\end{center}

\subsection{6. Summary: the connection is future
work}\label{summary-the-connection-is-future-work}

This supplement recorded a literature survey carried out after the
publication of the beat conjecture {[}1{]}.

\begin{itemize}
\tightlist
\item
  No prior work that proposes, tests, or observes the central mechanism
  of the beat conjecture (the emergence of Born statistics from the
  system--observer beat plus irrational-ratio Weyl equidistribution, a
  deviation in the coherent/commensurate regime) in the same form was
  found within the scope surveyed.
\item
  On the other hand, there exist prior programs that may share the
  problem-consciousness ``to derive the Born rule from a deterministic
  substrate + threshold detection as an emergent/approximate rule'' (La
  Cour--Williamson {[}3{]}, Khrennikov PCSFT {[}4,5{]}, the method of
  arbitrary functions {[}6{]}). The beat conjecture may be related to
  these, but whether it connects as a mechanism is future work.
\item
  A Born-rule test using a coherent LO (Conlon et al.~{[}7{]}) already
  exists. It is a null Sorkin higher-order interference test, and
  whether it connects to a measurement in the discriminating regime of
  this conjecture is future work.
\item
  Atomic clocks are the natural testbed of this conjecture, but the
  measurement needed for the discrimination (the commensurability
  dependence of single-shot projection statistics at a fixed relative
  phase) is, within the range surveyed, not performed; the prediction is
  not refuted but untested.
\end{itemize}

Whether these prior works connect logically to the beat conjecture, and
whether the discriminating measurement actually shows a deviation, are
both \textbf{future work}. This supplement cites only the necessary
parts where relevance is plausible and merely touches on the possible
relation. The status of the beat conjecture (a conjecture; the core
conditionally proven, the bridge unproven, single-shot-ness, the
exponent 2, and the scale attribution unsolved {[}1{]}) is unchanged by
this supplement.

\begin{center}\rule{0.5\linewidth}{0.5pt}\end{center}

\subsection{References}\label{references}

{[}1{]} N. Kihara, ``Emergence of Born Statistics from the
Equidistribution of the Observer--System Beat --- A Conjecture on the
Localized-Kernel Model,'' Zenodo, v0.4 (2026-06-28), Version DOI:
10.5281/zenodo.20985736 / Concept DOI: 10.5281/zenodo.20967081. (The
subject of this supplement = the preceding paper.)

{[}2{]} N. Kihara, ``Deriving the Form of the Born Distribution from the
Reproducing-Kernel Property of a Localized Odd-Harmonic Wave,'' Zenodo,
Concept DOI: 10.5281/zenodo.20965526.

{[}3{]} B. R. La Cour and M. C. Williamson, ``Emergence of the Born rule
in quantum optics,'' \emph{Quantum} \textbf{4}, 350 (2020). DOI:
10.22331/q-2020-10-26-350.

{[}4{]} A. Khrennikov, ``Detection model based on representation of
quantum particles by classical random fields: Born's rule and beyond,''
\emph{Phys. Lett. A} \textbf{372}, 6588--6592 (2008). DOI:
10.1016/j.physleta.2008.09.027.

{[}5{]} A. Khrennikov, ``Towards violation of Born's rule: description
of a simple experiment,'' \emph{AIP Conf. Proc.} \textbf{1327}, 387--393
(2011). DOI: 10.1063/1.3578726.

{[}6{]} L. Bonds, B. Burson, K. Cicchella, B. H. Feintzeig, Lynnx, A.
Yusaini, ``Quantum Probability via the Method of Arbitrary Functions,''
arXiv:2409.16457 (2024).

{[}7{]} L. O. Conlon, A. Walsh, Y. Hua, O. Thearle, T. Vogl, F.
Eilenberger, P. K. Lam, S. M. Assad, ``Testing the postulates of quantum
mechanics with coherent states of light and homodyne detection,''
\emph{New J. Phys.} \textbf{26}, 053003 (2024). DOI:
10.1088/1367-2630/ad3f1a.

{[}8{]} G. 't Hooft, \emph{The Cellular Automaton Interpretation of
Quantum Mechanics}, Fundamental Theories of Physics \textbf{185}
(Springer, 2016). DOI: 10.1007/978-3-319-41285-6.

{[}9{]} W. M. Itano, J. C. Bergquist, J. J. Bollinger, J. M. Gilligan,
D. J. Heinzen, F. L. Moore, M. G. Raizen, D. J. Wineland, ``Quantum
projection noise: Population fluctuations in two-level systems,''
\emph{Phys. Rev.~A} \textbf{47}, 3554 (1993). DOI:
10.1103/PhysRevA.47.3554.

{[}10{]} A. D. Ludlow, M. M. Boyd, J. Ye, E. Peik, P. O. Schmidt,
``Optical atomic clocks,'' \emph{Rev.~Mod. Phys.} \textbf{87}, 637
(2015). DOI: 10.1103/RevModPhys.87.637.

{[}11{]} M. Schioppo, R. C. Brown, W. F. McGrew, N. Hinkley, et al.,
``Ultrastable optical clock with two cold-atom ensembles,'' \emph{Nature
Photonics} \textbf{11}, 48--52 (2017). DOI: 10.1038/nphoton.2016.231.

{[}12{]} A. Al-Masoudi, S. Dörscher, S. Häfner, U. Sterr, C. Lisdat,
``Noise and instability of an optical lattice clock,'' \emph{Phys.
Rev.~A} \textbf{92}, 063814 (2015). DOI: 10.1103/PhysRevA.92.063814.

{[}13{]} B. K. Malia, J. Martínez-Rincón, Y. Wu, O. Hosten, M. A.
Kasevich, ``Free space Ramsey spectroscopy in rubidium with noise below
the quantum projection limit,'' \emph{Phys. Rev.~Lett.} \textbf{125},
043202 (2020). DOI: 10.1103/PhysRevLett.125.043202.

{[}14{]} M. Schulte, C. Lisdat, P. O. Schmidt, U. Sterr, K. Hammerer,
``Prospects and challenges for squeezing-enhanced optical atomic
clocks,'' \emph{Nat. Commun.} \textbf{11}, 5955 (2020). DOI:
10.1038/s41467-020-19403-7.

{[}15{]} H. Weyl, ``Über die Gleichverteilung von Zahlen mod. Eins,''
\emph{Math. Ann.} \textbf{77}, 313--352 (1916). DOI: 10.1007/BF01475864.

\begin{center}\rule{0.5\linewidth}{0.5pt}\end{center}

\emph{(This supplement records a literature survey carried out after the
publication of the preceding paper {[}1{]}. It claims neither identity
of mechanism nor experimental confirmation, and merely touches on the
possible relation within the range where relevance is plausible. The
logical connection to the preceding paper, and the execution of the
discriminating measurement, are left as future work.)}

\end{document}
